\section{Discussion}\label{sec:discussion}

\subsection{When Biological Structure Earns Complexity}

The two positive results share a common pattern: both work through
\emph{mechanism design} rather than \emph{algorithmic sophistication}.

The metabolic priority gate is a simple state machine.  When resources
drop below 10\%, operations with priority $< 5$ are rejected.  This is
not a complex algorithm---it is a structural rule.  But it provides a
hard guarantee: critical operations \emph{will} be served under any
load pattern, at the cost of bounded throughput reduction.  A flat
counter cannot provide this guarantee without reimplementing the state
machine, at which point it is no longer a flat counter.

The quorum sensing decay is exponential half-life applied to signal
concentrations.  Again, not algorithmically complex.  But the
combination of continuous signals (noise averaging) and temporal decay
(stale evidence removal) occupies a precision--recall position that
neither majority voting nor independent detection can reach without
adding similar structural properties.

In both cases, the biological metaphor provided the \emph{right
structural choice}---priority-dependent resource gating from metabolic
state machines, temporal signal decay from autoinducer kinetics---rather
than a better algorithm for the same structural choice.

\subsection{When It Does Not}

The epiplexity result demonstrates the failure mode of biological
design: adding algorithmic complexity (Bayesian two-signal detection)
over a simpler approach (cosine similarity) that already captures the
essential signal.

The theoretical appeal is sound.  Stagnation \emph{should} be
distinguishable from convergence by combining novelty and perplexity
signals.  But with mock embeddings, the novelty signal is essentially
random noise---two semantically identical messages with different
wording receive unrelated cosine similarity.  The second signal
(perplexity) is approximated heuristically when not directly measured.
With two noisy signals, the Bayesian combination performs worse than a
single clean signal.

This is consistent with the finding from Operon's C8 convergence
phase~\cite{banu2026operon}: biological abstractions generalize as
\emph{code structure} (the Genome round-trip is lossless; the
EpiplexityMonitor's distance provider works across abstraction scales)
but not as \emph{optimization algorithms} (tournament mutation matches
or outperforms LLM-guided evolution for configuration search).

The pattern: biological designs that enforce structural invariants
outperform simpler alternatives; biological designs that attempt
information processing do not reliably do so.

\subsection{Implications for Framework Integration}

Operon provides convergence adapters for six agent
frameworks~\cite{banu2026convergence}.  The benchmark results inform
which frameworks benefit most from biological structural guarantees:

\paragraph{A-Evolve.}
Its single-agent Solve-Observe-Evolve-Gate-Reload loop is where
metabolic budgets directly prevent unbounded exploration.  The
benchmark shows the fitness gate (a critical operation) is always
evaluated under metabolic budgeting---mapped to A-Evolve, this means
the evolutionary quality check is never skipped even when the
exploration budget is depleted.

\paragraph{Swarms.}
Its graph-based multi-agent topology is where quorum sensing provides
the most value: leaderless consensus with zero false positives.  When
false consensus triggers coordinated action across a swarm, the 0\%
FPR guarantee is operationally critical.

\paragraph{DeerFlow.}
Its hierarchical architecture with recursive sub-agent calls benefits
from metabolic budgets that prevent resource exhaustion.  The
MTORScaler's rate-of-change sensitivity is especially relevant for
DeerFlow's progressive skill loading, where resource consumption can
accelerate suddenly.

\paragraph{AnimaWorks, Ralph, Scion.}
Lighter benefits: genome immutability for configuration audit
(AnimaWorks), metabolic budgets mapping to backpressure enforcement
(Ralph), and quorum sensing for cross-container consensus (Scion).

\subsection{Relationship to Epistemic Topology}

Operon's epistemic topology layer derives four bounds from wiring
diagram structure: error amplification, sequential communication
overhead, parallel speedup, and tool density scaling.  These
operate at a different level than the mechanism benchmarks presented
here.

The epistemic bounds predict architecture-level properties (a long
pipeline \emph{will} accumulate sequential overhead); the mechanism
benchmarks test feature-level guarantees (critical operations
\emph{will} be served under pressure).  The two are complementary:
epistemic topology identifies architectures likely to need structural
guarantees; the mechanism benchmarks validate that those guarantees
work as claimed.

However, the epistemic bounds currently lack empirical validation
against measured behavior.  The metabolism and quorum sensing
benchmarks now have stronger empirical backing than the topology
theorems.  Validating the predicted error amplification and sequential
overhead bounds against actual benchmark measurements remains future
work.
